\section{The PAC/HAL/BSP Architecture Pattern}
\label{sec:pac-hal-bsp-architecture}

% Section 01.03: The PAC/HAL/BSP Architecture Pattern

Embedded Rust development follows a layered architecture that promotes code reuse, safety, and separation of concerns. This pattern is not arbitrary—it reflects years of lessons learned from implementing hardware support across multiple vendors and microcontroller families.

\subsection{The Three-Layer Architecture}
\label{subsec:three-layer-architecture}

\begin{figure}[H]
    \centering
    \begin{tikzpicture}[
        layer/.style={rectangle, draw=black, thick, minimum width=8cm, minimum height=1.5cm, font=\bfseries},
        arrow/.style={-Latex, thick, double},
        label/.style={font=\small\itshape}
    ]

        % Application Layer
        \node[layer, fill=rustblue!20] (app) at (0, 4.5) {Application Code};

        % BSP Layer
        \node[layer, fill=green!20] (bsp) at (0, 3) {Board Support Package};

        % HAL Layer
        \node[layer, fill=yellow!20] (hal) at (0, 1.5) {Hardware Abstraction Layer};

        % PAC Layer
        \node[layer, fill=red!20] (pac) at (0, 0) {Peripheral Access Crate};

        % Hardware Layer
        \node[layer, fill=gray!20] (hw) at (0, -1.5) {Hardware Peripherals};

        % Arrows showing dependency
        \draw[arrow] (app) -- (bsp);
        \draw[arrow] (bsp) -- (hal);
        \draw[arrow] (hal) -- (pac);
        \draw[arrow] (pac) -- (hw);

        % Labels
        \node[label, left=2cm of app] {Board-Specific};
        \node[label, left=2cm of bsp] {Portable};
        \node[label, left=2cm of hal] {Register Access};
        \node[label, left=2cm of pac] {Hardware};

        % Portability indicators
        \node[font=\small, right=2cm of app] {Application-Specific};
        \node[font=\small, right=2cm of bsp] {Board-Specific};
        \node[font=\small, right=2cm of hal] {Chip Family Specific};
        \node[font=\small, right=2cm of pac] {Vendor Provided / Generated};
        \node[font=\small, right=2cm of hw] {Fixed by Hardware};
    \end{tikzpicture}
    \caption{The PAC/HAL/BSP layered architecture}
    \label{fig:pac-hal-bsp-layers}
\end{figure}

Each layer has distinct responsibilities and portability characteristics.

\subsection{Peripheral Access Crate (PAC)}
\label{subsec:pac-layer}

The PAC provides the lowest-level, register-accurate representation of the microcontroller's peripherals. It is typically generated automatically from the vendor's System View Description (SVD) file using \codeinline{svd2rust}.

\textbf{Key characteristics}:
\begin{itemize}
    \item One-to-one mapping with hardware registers
    \item No abstractions on top of hardware
    \item All operations are \codeinline{unsafe} (by design)
    \item Generated automatically—never hand-edited
    \item Chip variant specific (EFR32MG24B220 vs A020)
\end{itemize}

\begin{minted}{rust}
// Example: Direct register access via PAC
unsafe {
    (*efr32mg24_pac::GPIO_S::ptr())
        .port_a_dout
        .modify(|r, w| w.bits(r.bits() | (1 << 5)));
}
\end{minted}

This writes directly to GPIO Port A, bit 5. The \codeinline{unsafe} keyword forces you to acknowledge that you're bypassing Rust's safety guarantees.

\begin{warningbox}
    PAC code should be wrapped in HAL abstractions. Direct PAC usage in application code is an anti-pattern that sacrifices safety for no benefit.
\end{warningbox}

\subsection{Hardware Abstraction Layer (HAL)}
\label{subsec:hal-layer}

The HAL provides safe, ergonomic APIs that wrap the raw PAC access. It implements standard traits for portability and adds compile-time safety guarantees.

\textbf{Key characteristics}:
\begin{itemize}
    \item Safe APIs (no \codeinline{unsafe} required by users)
    \item Type-safe state machines (e.g., pin modes encoded in types)
    \item Implements \crate{embedded-hal} traits
    \item Chip family specific (works across all EFR32MG24 variants)
    \item Medium abstraction level
    \item Written and maintained by community
\end{itemize}

\begin{minted}{rust}
// Example: GPIO via HAL
let pin = gpio.pin5.into_output();
pin.set_high();
\end{minted}

Compare this to the PAC version above. The HAL provides:
\begin{itemize}
    \item Type-safe pin selection (can't use non-existent pins)
    \item Compile-time mode tracking (can't read from output pin)
    \item Clear, readable method names
    \item No unsafe required
    \item Identical performance (zero-cost abstraction)
\end{itemize}

\subsection{Board Support Package (BSP)}
\label{subsec:bsp-layer}

The BSP provides board-specific abstractions and convenience APIs. It's the only layer that "knows" about your specific PCB design.

\textbf{Key characteristics}:
\begin{itemize}
    \item Maps logical functions to physical pins (e.g., LED1 → PA5)
    \item Provides board-specific initialization sequences
    \item Combines peripherals for board features
    \item Board specific (each PCB design has its own BSP)
    \item Highest abstraction level
    \item Minimal code (thin wrappers over HAL)
\end{itemize}

\begin{minted}{rust}
// Example: BSP for XIAO MG24 Sense
use xiao_mg24_bsp::{Board, LED};

let board = Board::new();
board.led_green.on();
delay_ms(500);
board.led_green.off();
\end{minted}

The BSP knows that the green LED connects to PA5 on this particular board and provides a convenient LED abstraction.

\subsection{The Flow of Abstractions}
\label{sec:abstraction-flow}

Here's how a simple operation (turning on an LED) flows through the layers:

\begin{enumerate}
    \item \textbf{Application}: Calls \codeinline{board.led_red.on()}
    \item \textbf{BSP}: Translates to \codeinline{hal.gpio.pin5.set_high()}
    \item \textbf{HAL}: Checks pin mode, then calls PAC
    \item \textbf{PAC}: Executes \codeinline{GPIO->DOUT |= (1 << 5)}
    \item \textbf{Hardware}: Physical pin voltage goes high, LED illuminates
\end{enumerate}

This layered approach provides benefits:

\textbf{Portability}: Switch to a different EFR32MG24 board? Only rewrite the BSP. HAL and PAC remain unchanged.

\textbf{Maintainability}: Bug in GPIO handling? Fix it in the HAL. All boards using that HAL benefit.

\textbf{Composability}: Libraries written for HAL traits work across all boards using that HAL.

\begin{notebox}
    \textbf{Architecture Rule}: Application code should never directly use PAC. Always go through HAL and BSP.
\end{notebox}

\subsection{Why This Separation Matters}
\label{sec:separation-importance}

The embedded Rust community settled on this architecture through experience:

\subsection{PAC Regeneration}

When vendors release updated SVD files (fixing errors or documenting new features), you regenerate the PAC without affecting HAL or BSP code.

\subsection{Multiple HAL Implementations}

Different use cases may need different HAL trade-offs:
\begin{itemize}
    \item High-performance HAL for real-time applications
    \item Small-footprint HAL for resource-constrained devices
    \item Feature-rich HAL for development boards
\end{itemize}

All can use the same PAC.

\subsection{Ecosystem Compatibility}

The \crate{embedded-hal} trait system enables an ecosystem of drivers:
\begin{itemize}
    \item Display drivers (SSD1306, ST7789) work with any SPI implementation
    \item Sensor drivers (BMP280, MPU6050) work with any I2C implementation
    \item No code changes needed when switching microcontrollers
\end{itemize}

\subsection{Our Approach in This Book}
\label{sec:our-approach}

In this book, we'll implement the stack bottom-up:

\begin{enumerate}
    \item \textbf{Chapter 2}: Generate the PAC from Silicon Labs SVD
    \item \textbf{Chapters 4-9}: Build HAL for all peripherals
    \item \textbf{Chapter 11}: RAIL radio abstraction
    \item \textbf{Chapters 12-14}: Protocol stacks (Bluetooth, Zigbee, Thread)
    \item \textbf{Chapter 15}: BSP for XIAO MG24 Sense
\end{enumerate}

This order ensures:
\begin{itemize}
    \item Each layer has a strong foundation
    \item Clear dependency management
    \item Comprehensive testing at each level
    \item Real hardware validation
\end{itemize}

\begin{exercise}[Easy]
    Look at a simple development board (like Arduino Uno) and list:
    \begin{enumerate}
        \item Which features belong in PAC vs HAL vs BSP?
        \item Why does Arduino's "digitalWrite()" combine all three layers?
    \end{enumerate}
\end{exercise}

\begin{exercise}[Medium]
    Design a type-safe state machine for a USART that:
    \begin{enumerate}
        \item Only allows baud rate configuration when USART is disabled
        \item Only allows transmission when USART is enabled
        \item Prevents mode changes during active transmission
    \end{enumerate}
    How would you encode these states in Rust's type system?
\end{exercise}

\begin{exercise}[Hard]
    Research other embedded HAL implementations:
    \begin{enumerate}
        \item Compare \codeinline{stm32-hal} vs \codeinline{nrf-hal}
        \item Note differences in how they handle:
            \begin{itemize}
                \item Clock configuration
                \item Error types
                \item Feature flags
                \item Documentation approach
            \end{itemize}
    \end{enumerate}
\end{exercise}
\clearpage
